\documentclass[11pt,a4paper]{article}

\usepackage[a4paper,margin=1in]{geometry}
\usepackage[T1]{fontenc}
\usepackage[utf8]{inputenc}
\usepackage{lmodern}
\usepackage{setspace}
\usepackage{parskip}
\usepackage{graphicx}
\usepackage{xcolor}
\usepackage{hyperref}
\usepackage{enumitem}
\usepackage{booktabs}

\hypersetup{
  colorlinks=true,
  linkcolor=black,
  urlcolor=blue,
  citecolor=black,
  pdftitle={CW3 SE Report Group 22}
}

\setstretch{1.08}
\setlist[itemize]{topsep=4pt,itemsep=2pt,leftmargin=1.6em}

\title{\textbf{CW3 SE Report Group 22}}
\author{}
\date{}

\begin{document}

\maketitle
\thispagestyle{empty}

\begin{center}
  Anonymous submission for the SE component of Coursework 3.
\end{center}

\vspace{1em}

% Keep this report anonymous:
% - do not include names
% - do not include UUNs
% - do not include the declaration of work here

\section{Task 1: Design Choices and Implementation Notes}

\subsection{Implemented Use Cases}

This submission implements the group-22 main-system use cases:
\begin{itemize}
  \item Log in
  \item Log out
  \item Register entertainment provider
  \item Create event
  \item Search for performances
  \item View performance
  \item Book performance
  \item Edit preferences
  \item Cancel performance
  \item Cancel booking
\end{itemize}

\subsection{System Structure}

Briefly describe the main architecture of the system. For example:
\begin{itemize}
  \item the main packages and their responsibilities
  \item the role of controllers, models, view, and external/integration components
  \item how data is currently stored during execution
\end{itemize}

\noindent
\textbf{Write here:}

\vspace{6em}

\subsection{Alignment with the Provided UML}

State whether the implementation matches the provided class and sequence diagrams exactly.
If there were any points where the diagrams could not be followed directly, explain:
\begin{itemize}
  \item what the mismatch was
  \item why it could not be implemented exactly as drawn
  \item what design change was made in code
  \item how the UML could be updated to stay consistent with the implementation
\end{itemize}

\noindent
\textbf{Write here:}

\vspace{8em}

\subsection{Key Design Choices}

Summarise the most important choices you made when implementing the main system. For example:
\begin{itemize}
  \item handling of input validation
  \item object ownership and shared state
  \item controller responsibilities
  \item interaction with mock external systems
  \item error handling strategy in the text interface
\end{itemize}

\noindent
\textbf{Write here:}

\vspace{8em}

\subsection{Quality Considerations in the Implementation}

Briefly mention implementation choices that support quality attributes such as security,
privacy, or reliability. This can overlap with the later reflection, but here the focus
should be on concrete system decisions.

\noindent
\textbf{Write here:}

\vspace{8em}

\section{Task 5: Reflection}

% Each of the following three parts should stay within 250 words.
% One or two paragraphs per part is enough.

\subsection{Reflection on Teamwork}

\noindent

\textbf{The Experience}
Our team first established a workflow using Git and GitHub with a fork-and-pull-request model.
We also built a CI check action for coursework. For Task 1, we began by dividing the use cases across the team before any implementation.
Once the use cases were divided, we wrote placeholder functions in all controller and model files.
This approach enabled a unified interface and facilitated assigning specific functions to each team member based on their use cases.

\textbf{Reflection on Action}
By using placeholder functions and assigning ownership from the start, we rarely encountered merge conflicts, as we seldom modified the same files or functions simultaneously.
This time, our engagement exceeded that of previous coursework.
We also implemented code review via pull requests, which was particularly important.
In a complex MVC system with shared data structures and interdependent controllers, this ensured teammates understood each other's APIs and data structures before integrating changes.
Our GitHub action CI check verified that CW3 files were modified, that all tests passed, and that the submission ZIP was built, which offers a convenient double-check after work.

\textbf{Theory}
We learned that defining function signatures early helps parallel development by reducing conflicts and forcing agreement on interfaces.
We also learned that code review is essential in complex systems, as it helps maintain a shared understanding across the team.

\textbf{Preparation}
Next time, we will adopt a Kanban-style task board to make the flow of work more visible and surface dependencies between teammates earlier,
rather than relying solely on ad-hoc messaging.

\subsection{Reflection on the Quality of the Work and Quality Attributes}

\noindent
\textbf{Word limit: 250 words maximum}

\textbf{Experience.}
We implemented all ten assigned use cases in a complex MVC system, then supported the
implementation with unit tests for the performance class and mock external systems, system tests for each use
case, and integrated tests that chain multiple use cases through shared state. We also
reviewed an external group's code submission as a structured code review exercise.

\textbf{Reflection on Action.}
The quality attribute best demonstrated by the final system is \textit{reliability} ---
the system behaves correctly and gracefully across a wide range of inputs. Every
controller method checks the current user's role before proceeding, preventing
unauthorised actions. All user inputs are parsed defensively: invalid types produce an
error and prompt a retry rather than crashing the system. Ownership is enforced before
any destructive operation --- only the student who created a booking can cancel it, and
only the entertainment provider who owns an event can cancel its performances. Payment
failures are handled gracefully: a booking is created first and marked
\texttt{PAYMENTFAILED} if the charge does not go through, preserving the booking record.
The cancellation workflow uses a two-pass approach --- refunds are processed for all
active bookings before any status is changed, so a failed refund aborts the entire
cancellation cleanly. These decisions are validated by 183 passing tests across unit,
system, and integrated levels.

\textbf{Theory.}
We learned that reliability emerges from many small, consistent guards rather than one
centralised check. Each null check, type parse, or role verification is simple in
isolation, but together they prevent the system from entering an inconsistent state.

\textbf{Preparation.}
With more time, we would address two gaps: passwords are currently stored in plaintext
in \texttt{docs/registered-eps.txt}, and in-memory data (events, performances, bookings)
is lost on restart. Password hashing and a persistent data store would be the natural
next steps.

\subsection{Reflection on Tools, Agile Approaches and Practices}

\textbf{Suggested structure:}
\begin{itemize}
  \item why each tool, practice, or agile approach was chosen
  \item what advantages and disadvantages were observed
  \item what worked well in practice and what did not
  \item what would be changed next time
\end{itemize}

\noindent
\textbf{Word limit: 250 words maximum}

\vspace{2em}

\subsection{Screenshots of Tools and Practices}

The coursework document explicitly asks for screenshots in the tools/agile/practices reflection.
Insert screenshots of the tools you actually used, such as issue tracking, Kanban boards,
version control history, CI runs, or planning boards.

\begin{figure}[h]
  \centering
  % Replace the filename below with your real screenshot.
  \fbox{\rule{0pt}{2in}\rule{0.9\linewidth}{0pt}}
  \caption{Screenshot placeholder 1. Replace with a real screenshot of a tool or practice used by the team.}
\end{figure}

\begin{figure}[h]
  \centering
  % Replace the filename below with your real screenshot.
  \fbox{\rule{0pt}{2in}\rule{0.9\linewidth}{0pt}}
  \caption{Screenshot placeholder 2. Replace with another real screenshot if needed.}
\end{figure}

\section{Submission Checklist}

Before exporting the PDF, check the following:
\begin{itemize}
  \item The file is named \texttt{CW3SEReportGroup22.pdf}.
  \item The document is anonymous.
  \item The declaration of work is not included here.
  \item Task 1 choices are described.
  \item All three reflection parts are included.
  \item Each reflection part is at most 250 words.
  \item The tools/agile section includes screenshots.
\end{itemize}

\end{document}
